
% Szkielet dla pracy pisanej w języku polskim.

\documentclass[polish,a4paper,oneside]{ppfcmthesis}


\usepackage[utf8]{inputenc}
\usepackage[OT4]{fontenc}


\author{Ignacy Iksiński}                              % Your name comes here
\title{W zdrowym ciele zdrowy~duch}                   % Note how we protect the final title phrase from breaking
\ppsupervisor{prof.~dr hab.~inż.~Alojzy Wołodyjowski} % Your supervisor comes here.
\ppyear{2006}                                         % Year of final submission (not graduation!)


\begin{document}

% Front matter starts here
\frontmatter\pagestyle{empty}%
\maketitle\cleardoublepage%

% Blank info page for "karta dyplomowa"
\thispagestyle{empty}\vspace*{\fill}%
\begin{center}Tutaj przychodzi karta pracy dyplomowej;\\oryginał wstawiamy do wersji dla archiwum PP, w pozostałych kopiach wstawiamy ksero.\end{center}%
\vfill\cleardoublepage%

% Table of contents.
\pagenumbering{Roman}\pagestyle{ppfcmthesis}%
\tableofcontents* \cleardoublepage%

% Main content of your thesis starts here.
\mainmatter%

\chapter{Wstęp}

Wstęp\footnote{Treść przykładowych rozdziałów została skopiowana
z ,,zasad'' redakcji prac dyplomowych FCMu~\cite{fcm-red}.} do pracy powinien zawierać następujące elementy:
\begin{itemize}
    \item krótkie uzasadnienie podjęcia tematu; 
    \item cel pracy (patrz niżej); 
    \item zakres (przedmiotowy, podmiotowy, czasowy) wyjaśniający, w jakim rozmiarze praca będzie realizowana; 
    \item ewentualne hipotezy, które autor zamierza sprawdzić lub udowodnić; 
    \item krótką charakterystykę źródeł, zwłaszcza literaturowych; 
    \item układ pracy (patrz niżej), czyli zwięzłą charakterystykę zawartości poszczególnych rozdziałów; 
    \item ewentualne uwagi dotyczące realizacji tematu pracy np.~trudności, które pojawiły się w trakcie 
    realizacji poszczególnych zadań, uwagi dotyczące wykorzystywanego sprzętu, współpraca z firmami zewnętrznymi. 
\end{itemize}

\noindent
\textbf{Wstęp do pracy musi się kończyć dwoma następującymi akapitami:}
\begin{quote}
Celem pracy jest opracowanie / wykonanie analizy / zaprojektowanie / ...........
\end{quote}
oraz:
\begin{quote}
Struktura pracy jest następująca. W rozdziale 2 przedstawiono przegląd literatury na temat ........ 
Rozdział 3 jest poświęcony ....... (kilka zdań). 
Rozdział 4 zawiera ..... (kilka zdań) ............ itd. 
Rozdział X stanowi podsumowanie pracy. 
\end{quote}

W przypadku prac inżynierskich zespołowych lub magisterskich 2-osobowych, po tych dwóch w/w akapitach 
musi w pracy znaleźć się akapit, w którym będzie opisany udział w pracy poszczególnych członków zespołu. Na przykład:

\begin{quote}
Jan Kowalski w ramach niniejszej pracy wykonał projekt tego i tego, opracował ......
Grzegorz Brzęczyszczykiewicz wykonał ......, itd. 
\end{quote}



\chapter{Opis procesów biznesowych}
\section{Aktorzy}
\section{Obiekty biznesowe}
\section{Biznesowe przypadki użycia}

Rozdział teoretyczny --- przegląd literatury naświetlający stan wiedzy na dany temat. 

Przegląd literatury naświetlający stan wiedzy na dany temat obejmuje rozdziały pisane na podstawie
literatury, której wykaz zamieszczany jest w części pracy pt.~\emph{Literatura} (lub inaczej \emph{Bibliografia},
\emph{Piśmiennictwo}). W tekście pracy muszą wystąpić odwołania do wszystkich pozycji zamieszczonych w
wykazie literatury. \textbf{Nie należy odnośników do literatury umieszczać w stopce strony.} Student jest
bezwzględnie zobowiązany do wskazywania źródeł pochodzenia informacji przedstawianych w pracy,
dotyczy to również rysunków, tabel, fragmentów kodu źródłowego programów itd. Należy także podać
adresy stron internetowych w przypadku źródeł pochodzących z Internetu.

\chapter{Wymagania funkcjonalne}
\section{Wstęp - diagram przypadków użycia}
\section{Administrator iQuest}
\subsection{UC1}
\subsection{UC2}
\section{Ankieter iQuest}
\subsection{UC3}
\subsection{UC4}
\section{Respondent iQuest}
\subsection{UC5}
\subsection{UC6}

\chapter{Wymagania pozafunkcjonalne}
\section{Wstęp}
\section{Charakterystyki oprogramowania}
\section{Wymagania pozafunkcjonalne i ich weryfikacja}

\chapter{Architektura systemu}
\section{Wstęp}
\section{Opis ogólny architektury}
\subsection{Perspektywa fizyczna}
\subsection{Perspektywa implementacyjna}
\section{Komunikacja pomiędzy modułami}
\section{Schemat bazy danych}
\section{Perspektywa kodu}

\chapter{}
\chapter{}
\chapter{}
\chapter{}
\chapter{}


\chapter{Rozwinięcie}

Rozdziały dokumentujące pracę własną studenta: opisujące ideę, sposób lub metodę 
rozwiązania postawionego problemu oraz rozdziały opisujące techniczną stronę rozwiązania 
--- dokumentacja techniczna, przeprowadzone testy, badania i uzyskane wyniki. 

Praca musi zawierać elementy pracy własnej autora adekwatne do jego wiedzy praktycznej uzyskanej w
okresie studiów. Za pracę własną autora można uznać np.: stworzenie aplikacji informatycznej lub jej
fragmentu, zaproponowanie algorytmu rozwiązania problemu szczegółowego, przedstawienie projektu 
np.~systemu informatycznego lub sieci komputerowej, analizę i ocenę nowych technologii lub rozwiązań
informatycznych wykorzystywanych w przedsiębiorstwach, itp. 

Autor powinien zadbać o właściwą dokumentację pracy własnej obejmującą specyfikację założeń i 
sposób realizacji poszczególnych zadań
wraz z ich oceną i opisem napotkanych problemów. W przypadku prac o charakterze 
projektowo-implementacyjnym, ta część pracy jest zastępowana dokumentacją techniczną i użytkową systemu. 

W pracy \textbf{nie należy zamieszczać całego kodu źródłowego} opracowanych programów. Kod źródłowy napisanych
programów, wszelkie oprogramowanie wytworzone i wykorzystane w pracy, wyniki przeprowadzonych
eksperymentów powinny być umieszczone na płycie CD, stanowiącej dodatek do pracy.

\section*{Styl tekstu}

Należy\footnote{Uwagi o stylu pochodzą częściowo ze stron Macieja Drozdowskiego~\cite{mdro}.} 
stosować formę bezosobową, tj.~\emph{w pracy rozważono ......, 
w ramach pracy zaprojektowano ....}, a nie: \emph{w pracy rozważyłem, w ramach pracy zaprojektowałem}. 
Odwołania do wcześniejszych fragmentów tekstu powinny mieć następującą postać: ,,Jak wspomniano wcześniej, ....'', 
,,Jak wykazano powyżej ....''. Należy unikać długich zdań. 

,,Ilość'' i ,,liczba''. Proszę zauważyć, liczba dotyczy rzeczy policzalnych, np.~liczba osób, liczba zadań, procesorów. 
Ilość dotyczy rzeczy niepoliczalnych, np.~ilość wody, energii. Należy starać się wyrażać precyzyjnie, tj.~zgodnie 
z naturą liczonych obiektów.\footnote{(DW) Według wytycznych Rady Języka Polskiego obie formy są dopuszczalne
zarówno do obiektów policzalnych, jak i niepoliczalnych. W tekstach technicznych warto być jednak precyzyjnym.}

Niedopuszczalne są zwroty używane w języku potocznym. W pracy należy używać terminologii informatycznej, która ma 
sprecyzowaną treść i znaczenie. Nie należy używać ,,gazetowych'' określeń typu: 
silnik bazy danych, silnik programu, maszyna skryptowa, elektroniczny mechanizm, mapowanie, string, gdyż nie wiadomo 
co one właściwie oznaczają. 

Niedopuszczalne jest pisanie pracy metodą \emph{cut\&paste}, bo jest to plagiat i dowód intelektualnej indolencji autora.
Dane zagadnienie należy opisać własnymi słowami. Zawsze trzeba powołać się na zewnętrzne źródła. 



\chapter{Zakończenie}

Zakończenie pracy zwane również Uwagami końcowymi lub Podsumowaniem powinno zawierać ustosunkowanie
się autora do zadań wskazanych we wstępie do pracy, a w szczególności do celu i zakresu pracy oraz
porównanie ich z faktycznymi wynikami pracy. Podejście takie umożliwia jasne określenie stopnia
realizacji założonych celów oraz zwrócenie uwagi na wyniki osiągnięte przez autora w ramach jego
samodzielnej pracy.

Integralną częścią pracy są również dodatki, aneksy i załączniki np.~płyty CDROM
zawierające stworzone w ramach pracy programy, aplikacje i projekty.


% All appendices and extra material, if you have any.
\cleardoublepage\appendix%

\chapter{Parę słów o stylu \texttt{ppfcmthesis}}

\section{Różnice w stosunku do ,,oficjalnych'' zasad składu ze stron FCMu}

Autor niniejszego stylu nie zgadza się z niektórymi zasadami wprowadzonymi w oficjalnym 
dokumencie FCMu.\footnote{\url{http://www.fcm.put.poznan.pl/platon/dokumenty/dlaStudentow/egzaminDyplomowy/zasadyRedakcji}}
Poniższe elementy są składane nieco inaczej w stosunku do ,,oficjalnych'' wytycznych.

\begin{itemize}
    \item Promotor na stronie tytułowej jest umiejscowiony w centralnej osi pionowej strony (a
    nie po prawej stronie).
    
    \item Czcionka użyta do składu to nie Times New Roman.
    
    \item Spacje między tytułami akapitów oraz wcięcia zostały pozostawione takie, jak są zdefiniowane
    oryginalnie w pakiecie Memoir (oraz w \LaTeX{}u). Jeśli zdefiniowano ,,polską'' opcję składu,
    to będzie w użyciu wcięcie pierwszego akapitu po tytułach rozdziałów. Przy składzie ,,angielskim''
    tego wcięcia nie ma.

    \item Odwrócona jest kolejność rozdziałów \emph{Literatura} i \emph{Dodatki}.

    \item Na ostatniej stronie umieszczono stopkę informującą o prawach autorskich i programie
    użytym do składu.
    
    \item Nie do końca zgadzam się ze stwierdzeniem, iż ,,zamieszczanie list tabel, rysunków, 
    wykresów w pracy dyplomowej jest nieuzasadnione''. Niektóre typy publikacji zawierają tabele i rysunki, których
    skorowidz umożliwia łatwiejsze ich odszukanie. Ale niech będzie.

    \item Styl podpisów tabel jest taki sam, jak rysunków i odmienny od FCMowego. 
    Jeśli ktoś koniecznie chce mieć zgodne z wytycznymi
    podpisy, to zamiast \texttt{caption} niech użyje \texttt{fcmtcaption} do podpisywania tablic oraz
    \texttt{fcmfcaption} do podpisywania rysunków. Podpisy pod rysunkami pozostaną pełne, a nie skrócone (,,Rys.'').
    
    \item Styl formatowania literatury jest nieco inny niż proponowany przez FCM.
\end{itemize}



\chapter{Składanie dokumentu w systemie \LaTeX}

Po pierwsze to gratulacje --- dobry wybór. W tym rozdziale znajduje się
garść informacji o tym, jak poprawnie składać tekst pracy w systemie \LaTeX{} wraz z 
przykładami, które mają służyć do przeklejania do własnych dokumentów.

\section{Narzędzia}
Pracując pod systemem Windows, polecam:
\begin{itemize}
    \item MikTeX, \url{http://www.miktex.org/},
    \item JEdit, \url{http://www.jedit.org/},
    \item Ghostview, Ghostscript (podgląd dokumentów PDF bez blokowania pliku):\\
        \url{http://www.cs.wisc.edu/~ghost/}. 
\end{itemize}

Po zainstalowaniu tych narzędzi wystarczy wykonać polecenie \texttt{compile.bat} (który
jest skryptem wsadowym dla Windows). Dla tych, którzy wolą nieco automatyzacji --- skrypt
\texttt{latexmk}, który jest w MikTeXu (a który potrzebuje zainstalowanego Perla) jest
również bardzo wygodny: \texttt{latexmk -pdf -pvc main.tex}.

\section{Edycja tekstu}
\chaptermark{Tytuł rozdziału, jeśli pełen się nie mieści\ldots{}}{}

\subsection{Struktura dokumentu}

Praca składa się z rozdziałów (\texttt{chapter}) i podrozdziałów (\texttt{section}).
Ewentualnie można również rozdziały zagnieżdzać (\texttt{subsection}, \texttt{subsubsection}),
jednak nie powinno się wykraczać poza drugi poziom hierarchii (czyli \texttt{subsubsection}).

\subsection{Akapity i znaki specjalne}

Każdy akapit to po prostu blok tekstu. Nieważne jak sformatowany --- to zrobi już
system $\LaTeX$.

Akapity rozdziela się od siebie przynajmniej jedną pustą linią. Podstawowe
instrukcje, które się przydają to \emph{wyróżnienie pewnych słów}. Można również
stosować \textbf{styl pogrubiony}, choć nie jest to generalnie zalecane.

Należy pamiętać o zasadach polskiej interpunkcji i ortografii. Po spójnikach 
jednoliterowych warto wstawić znak tyldy ($\sim$), który jest tak zwaną
,,twardą spacją'' i powoduje, że wyrazy nią połączone nie będą rozdzielane
na dwie linie tekstu.

Polskie znaki interpunkcyjne różnią się nieco od angielskich: to jest ,,polski'', a to jest
``angielski''. W kodzie źródłowym tego tekstu będzie widać różnicę.

Proszę również zwrócić uwagę na znak myślnika, który może być pauzą ,,---'' lub
półpauzą: ,,--''. Należy stosować je konsekwentnie. Do łączenia wyrazów używamy
zwykłego ,,-'' (\emph{północno-wschodni}), do myślników --- pauzy lub półpauzy.
Inne zasady interpunkcji i typografii można znaleźć w słownikach.

\subsection{Wypunktowania}

Wypunktowanie z cyframi:
\begin{enumerate}
    \item to jest punkt,
    \item i to jest punkt,
    \item a to jest ostatni punkt.
\end{enumerate}

\noindent
Po wypunktowaniach czasem nie warto wstawiać wcięcia akapitowego. Wtedy przydatne jest
polecenie \texttt{noindent}. Wypunktowanie z kropkami (tzw.~\emph{bullet list}) wygląda tak:
\begin{itemize}
    \item to jest punkt,
    \item i to jest punkt,
    \item a to jest ostatni punkt.
\end{itemize}

\noindent
Wypunktowania opisowe właściwie niewiele się różnią:
\begin{description}
    \item[elementA] to jest opis,
    \item[elementB] i to jest opis,
    \item[elementC] a to jest ostatni opis.
\end{description}


\subsection{Polecenia pakietu \texttt{ppfcmthesis}}

Parę poleceń zostało zdefiniowanych aby uspójnić styl pracy. Są one przedstawione poniżej
(oczywiście nie trzeba się do nich stosować).

\paragraph{Makra zdefiniowane dla języka angielskiego.} Są nimi: \texttt{termdef} oraz \texttt{acronym}.
Przykłady poniżej obrazują ich przewidywane użycie w tekście.
\begin{center}\footnotesize%
\begin{tabular}{l >{\rightskip\fill}p{12cm}}
\toprule
źródło   & \texttt{we call this a $\backslash$termdef\{Database Management System\} ($\backslash$acronym\{DBMS\})} \\ \cmidrule(lr){2-2}
docelowo & we call this a \termdef{Database Management System} (\acronym{DBMS}) \\ 
\bottomrule
\end{tabular}
\end{center}

\paragraph{Makra zdefiniowane dla języka polskiego.} Podobnie jak dla języka angielskiego zdefiniowano
odpowiedniki polskie: \texttt{defini\-cja}, \texttt{akronim} oraz \texttt{english} dla tłumaczeń angielskich
terminów. Przykłady poniżej obrazują ich przewidywane użycie w tekście.
\begin{center}\footnotesize%
\begin{tabular}{l >{\rightskip\fill}p{12cm}}
\toprule
źródło   & \texttt{nazywamy go $\backslash$definicja\{systemem zarządzania bazą danych\} ($\backslash$akronim\{DBMS\}, $\backslash$english\{Database Management System\})} \\ \cmidrule(lr){2-2}
docelowo & nazywamy go \definicja{systemem zarządzania bazą danych} (\akronim{DBMS}, \english{Database Management System}) \\ \bottomrule
\end{tabular}
\end{center}


\subsection{Rysunki}

Format wstawianych rysunków zależy od tego czy używa się do kompilacji polecenia
\texttt{latex}, czy też \texttt{pdflatex}. Oba powinny dać dokładnie ten sam wynik końcowy,
ale praca z nimi jest nieco inna.

\begin{description}
    \item[latex] To polecenie kompiluje źródła \LaTeX{}owe do pliku 
        z rozszerzeniem \texttt{dvi}. Ten plik można przeglądać przy pomocy specjalizowanych programów
        takich jak przykładowo Yap obecny z dystrybucją Mik\TeX{}a. Aby uzyskać docelowy plik \akronim{PDF}
        należy przekonwertować plik \texttt{dvi} przy pomocy programu \texttt{dvipdfm}. 
        
        \textbf{UWAGA:} korzystając z programu \texttt{latex}, wszystkie rysunki muszą być w formacie \akronim{EPS}
        (\english{encapsulated postscript}).

    \item[pdflatex] To polecenie kompiluje źródła \LaTeX{}owe bezpośrednio do pliku \akronim{PDF}.
    
        \textbf{UWAGA:} korzystając z programu \texttt{pdflatex}, wszystkie rysunki muszą być w formacie \akronim{PDF},
        \akronim{JPG} lub \akronim{PNG}.
\end{description}

Można oczywiście używać obu systemów --- wtedy pliki rysunków muszą po prostu być dostępne w obu formatach.

Wszystkie rysunki (w tym również diagramy, szkice i inne) osadzamy w środowisku 
\texttt{figure} i umieszczamy podpis \emph{pod} rysunkiem, w formie elementu \texttt{caption}. Rysunki powinny
zostać umieszczone u góry strony (osadzone bezpośrednio w treści strony zwykle utrudniają czytanie tekstu).
Rysunek~\ref{rys:plama} zawiera przykład pełnego osadzenia rysunku na stronie.

\begin{figure}[t] % możliwe opcje to 't' - top, 'b' - bottom, 'h' - 'here', ale zaleca się 't'
\centering\includegraphics[width=5cm]{figures/template/logo-pp}
\caption{Logo Politechniki Poznańskiej.}\label{rys:plama}
\end{figure}

\begin{figure}[t]
\centering\includegraphics[width=5cm]{figures/template/logo-pp}
\fcmfcaption{Logo Politechniki Poznańskiej. Formatowanie zgodne z wytycznymi FCMu.}\label{rys:plama2}
\end{figure}

Zasady FCMu sugerują nieco inne nagłówki rysunków. Dostepne są one poleceniem \texttt{fcmfcaption} (zob.~rysunek
\ref{rys:plama2}), jeśli ktoś woli mieć podpisy niespójne z rysunkami\ldots

\subsection{Tablice}

Tablice to piękna rzecz, choć akurat ich umiejętne tworzenie w \LaTeX{}u nie jest łatwe. 
Jeśli tablica jest skomplikowana, to pewnie łatwiej będzie ją wykonać w programie
OpenOffice, a następnie wyeksportować jako plik \akronim{PDF}. W każdym przypadku tablice wstawia się podobnie
jak rysunki, tylko że w środowisko \texttt{table}. Tradycja typograficzna sugeruje umieszczenie opisu tablicy, a więc
elementu \texttt{caption} ponad jej treścią (inaczej niż przy rysunkach).  

Tablica~\ref{tab:tabela} pokazuje pełen przykład.

\begin{table}[h]
\caption{Przykładowa tabela. Styl opisu jest zgodny z rysunkami.}\label{tab:tabela}
\centering\footnotesize%
\begin{tabular}{l c}
\toprule
artykuł & cena [zł] \\
\midrule
bułka   & $0,4$ \\
masło   & $2,5$ \\
\bottomrule
\end{tabular}
\end{table}

Zasady FCMu sugerują nieco inne nagłówki tablic. Dostepne są one poleceniem \texttt{fcmtcaption} (zob.~tablicę
\ref{tab:tabela2}), jeśli ktoś woli mieć podpisy niespójne z rysunkami\ldots

\begin{table}[h]
\fcmtcaption{Przykładowa tabela. Styl opisu jest zgodny z wytycznymi FCMu.}\label{tab:tabela2}
\centering\footnotesize%
\begin{tabular}{l c}
\toprule
artykuł & cena [zł] \\
\midrule
bułka   & $0,4$ \\
masło   & $2,5$ \\
\bottomrule
\end{tabular}
\end{table}


\section{Literatura i materiały dodatkowe}

Materiałów jest mnóstwo. Oto parę z nich:
\begin{itemize}
    \item \emph{The Not So Short Introduction\ldots}, która posiada również tłumaczenie 
    w języku polskim.\\
    \url{http://www.ctan.org/tex-archive/info/lshort/english/lshort.pdf}

    \item Klasy stylu \texttt{memoir} posiadają bardzo wiele informacji o składzie tekstów
    anglosaskich oraz sposoby dostosowania \LaTeX{}a do własnych potrzeb.\\
    \url{http://www.ctan.org/tex-archive/macros/latex/contrib/memoir/memman.pdf}
    
    \item Nasza grupa dyskusyjna i repozytorium SVN są również dobrym miejscem aby zapytać
    (lub sprawdzić czy pytanie nie zostało już zadane).\\
    \url{https://ophelia.cs.put.poznan.pl/svn/put-latex/trunk}

    \item Dla łaknących więcej wiedzy o systemie \LaTeX{} podstawowym źródłem informacji
    jest książka Lamporta~\cite{Lamport:LDP85}. Prawdziwy \emph{hardcore} to oczywiście
    \emph{The \TeX{}book} profesora Knutha~\cite{Knuth:ct-a}.
\end{itemize}



% Bibliography (books, articles) starts here.
\bibliographystyle{plalpha}{\raggedright\sloppy\small\bibliography{bibliography}}

% Colophon is a place where you should let others know about copyrights etc.
\ppcolophon

\end{document}
